Here, to verify whether individual highly cited papers have a longterm impact, we analyze their lifespan structure. First, in order to determine whether a paper has reached an obsolescence state, we adopt the longterm citation model of P.D.B. Parolo etal.\cite{Par1} and the citationdecay model of D.Wang etal.\cite{Wan1}, and set the following three conditions: (1) a sufficient period has elapsed since the publication year, (2) the peak period has sufficiently passed, and (3) the citations have sufficiently decayed. The specific threshold values for these three conditions are shown in Figure\ref{fig:gr6}. As a result of this assessment, 77 out of 536 papers were judged to be in an obsolescence state. The indicator values for the highly cited papers judged to be in an obsolescence state and for all highly cited papers are shown in Figure\ref{fig:arr1}. For the papers judged to be in an obsolescence state, the mode of the period from publication to peak is 4years, the mode of the period from peak to sufficient citation decay is 7years, and the mode of the period from publication to sufficient citation decay is 12years. This result supports the report by P.D.B. Parolo etal. (in the lifescience field, the peak occurs within 5years after publication). In contrast, for all highly cited papers the period from publication to peak is around 10years, showing a different pattern. A plot of citation counts for the highly cited papers judged to be in an obsolescence state is shown in Figure\ref{fig:gr7}. The period from publication to sufficient citation decay ranges from a minimum of 10years to a maximum of 84years. From this figure, it can be inferred that more recent papers tend to receive high citations immediately after publication.
